\chapter{Остатки, сравнения и целочисленные уравнения}\markright{Глава 21}

\section*{Что сохраняется, когда полные обороты отброшены}
На часах после двенадцати снова появляется один, а дни недели повторяются через семь дней. Вместо огромного числа удобно хранить его остаток. Это не приближённое вычисление: при фиксированном модуле остаток точно описывает, какие числа отличаются на целое число полных циклов.

\begin{definition}
\term{Деление с остатком} (division with remainder) целого \(a\) на натуральное \(m>0\) — представление \(a=mq+r\), где \(q\in\mathbb Z\) — \term{неполное частное} (quotient), а \(r\in\mathbb Z\), \(0\le r<m\), — \term{остаток} (remainder).
\end{definition}
Например, \(47=5\cdot9+2\), а \(-47=5\cdot(-10)+3\). Остаток по принятому определению не бывает отрицательным. Поэтому \(-47=5\cdot(-9)-2\) не является делением с остатком.

\begin{theorem}[Существование и единственность остатка]
При любом целом \(a\) и натуральном \(m>0\) существуют единственные целые \(q,r\), для которых \(a=mq+r\) и \(0\le r<m\).
\end{theorem}
\par\addvspace{5pt}\begin{lrbox}{\CatFlowBox}\begin{minipage}[t]{\linewidth}
\setlength{\parskip}{5pt}\setlength{\columnsep}{4.5mm}\setlength{\intextsep}{0pt}
\begin{wrapfigure}[13]{r}{54mm}
\includegraphics[width=54mm,height=58mm,keepaspectratio]{images/unique/art-59.png}
\end{wrapfigure}
\textit{Доказательство.} Среди чисел \(a-mq\), где \(q\in\mathbb Z\), есть неотрицательные: достаточно взять \(q=-|a|\). Выберем среди них наименьшее \(r\). Здесь используется принцип наименьшего элемента: непустое подмножество \(\mathbb N_0\) имеет наименьший элемент. Он следует из индукции: если бы множество не имело наименьшего элемента, последовательно доказывалось бы, что в нём нет ни \(0\), ни \(1\), ни любого следующего числа. Если \(r\ge m\), то \(r-m=a-m(q+1)\) — меньший неотрицательный кандидат, противоречие. Значит, \(r<m\).
Если также \(a=mq'+r'\), то \(m(q-q')=r'-r\). Правая часть по модулю меньше \(m\); единственное кратное \(m\) с таким свойством — ноль. Поэтому \(r=r'\) и \(q=q'\).\hfill\(\blacksquare\)
\par\WFclear\end{minipage}\end{lrbox}%
\ifdim\dimexpr\ht\CatFlowBox+\dp\CatFlowBox\relax<64mm\dp\CatFlowBox=\dimexpr64mm-\ht\CatFlowBox\relax\fi%
\noindent\usebox{\CatFlowBox}\par\addvspace{5pt}


\section*{Сравнения по модулю}
\begin{definition}
Целые \(a,b\) \term{сравнимы по модулю} (congruent modulo) натурального \(m>0\), если \(m\mid(a-b)\). Запись: \(a\equiv b\pmod m\). Число \(m\) называют \term{модулем сравнения} (modulus).
\end{definition}
Сравнимость означает совпадение остатков. Например, \(22\equiv4\pmod6\), \(-2\equiv4\pmod6\). Запись \(22\equiv4\pmod6\) не утверждает обычное равенство \(22=4\). Модуль должен быть указан: \(22\) и \(4\) сравнимы по модулю \(6\), но не по модулю \(5\).

\begin{theorem}[Арифметика сравнений]
Если \(a\equiv b\pmod m\) и \(c\equiv d\pmod m\), то
\[
a+c\equiv b+d\pmod m,\qquad a-c\equiv b-d\pmod m,\qquad ac\equiv bd\pmod m.
\]
В частности, \(a^n\equiv b^n\pmod m\) для любого \(n\in\mathbb N\).
\end{theorem}
\begin{proof}
Пусть \(a-b=mu\), \(c-d=mv\) с целыми \(u,v\). Разность сумм равна \(m(u+v)\), разность разностей — \(m(u-v)\). Для произведения
\[
ac-bd=(a-b)c+b(c-d)=m(uc+bv).
\]
Все разности делятся на \(m\). Утверждение о натуральных степенях следует последовательным применением правила произведения, то есть индукцией по \(n\).
\end{proof}
Пример: \(10\equiv1\pmod9\), поэтому \(10^{100}+2\equiv3\pmod9\). Для вычисления последней цифры выбирают модуль \(10\); для дня недели — \(7\). До возведения в большую степень полезно выписать несколько первых остатков и заметить цикл, затем обосновать его повторение умножением на одно и то же число.

\begin{definition}
\term{Класс вычетов} (residue class) числа \(a\) по модулю \(m\) — множество \([a]_m=\{a+mk:k\in\mathbb Z\}\). Множество всех классов обозначают \(\mathbb Z/m\mathbb Z\). При \(m\ge2\) оно состоит ровно из \(m\) классов \([0]_m,\ldots,[m-1]_m\).
\end{definition}
Числа \(a\) и \(a+m\) различны, но задают один класс. По предыдущей теореме правила \([a]_m+[b]_m=[a+b]_m\) и \([a]_m[b]_m=[ab]_m\) не зависят от выбора представителей. Это первый пример нового множества, на котором можно определить знакомые действия. В следующей главе мы сравним его свойства со свойствами целых и рациональных чисел.

\section*{Тождество Безу и обратный проход алгоритма Евклида}
\begin{definition}
\term{Целочисленная линейная комбинация} (integer linear combination) целых \(a,b\) — число вида \(au+bv\), где \(u,v\in\mathbb Z\). Равенство \(\gcd(a,b)=au+bv\) называют \term{тождеством Безу} (Bézout's identity). Это другое утверждение, чем теорема Безу о корне многочлена из главы 11.
\end{definition}

\begin{theorem}[Безу для целых чисел]
Если \(a,b\) — целые, не равные нулю одновременно, то существуют целые \(u,v\) с \(au+bv=\gcd(|a|,|b|)\). Здесь \(\gcd(n,0)=n\) при \(n>0\).
\end{theorem}
\par\addvspace{5pt}\begin{lrbox}{\CatFlowBox}\begin{minipage}[t]{\linewidth}
\setlength{\parskip}{5pt}\setlength{\columnsep}{4.5mm}\setlength{\intextsep}{0pt}
\begin{wrapfigure}[10]{l}{54mm}
\includegraphics[width=54mm,height=58mm,keepaspectratio]{images/unique/art-60.png}
\end{wrapfigure}
\textit{Доказательство.} Для положительных \(|a|,|b|\) алгоритм Евклида последовательно делит с остатком. Каждый новый остаток равен предыдущему делимому минус целое кратное делителя. Поэтому первый остаток выражается через \(|a|,|b|\), следующий — тоже, и по индукции каждый остаток является их целочисленной линейной комбинацией. Последний ненулевой остаток — НОД, поскольку при замене пары \((A,B)\) на \((B,A-qB)\) множество общих делителей не меняется. Знаки исходных \(a,b\) учитываются изменением знаков коэффициентов. Если одно число равно нулю, достаточно взять коэффициент при другом равным \(1\) или \(-1\).\hfill\(\blacksquare\)
\par\WFclear\end{minipage}\end{lrbox}%
\ifdim\dimexpr\ht\CatFlowBox+\dp\CatFlowBox\relax<64mm\dp\CatFlowBox=\dimexpr64mm-\ht\CatFlowBox\relax\fi%
\noindent\usebox{\CatFlowBox}\par\addvspace{5pt}


Пример обратной подстановки:
\[
252=198+54,\quad198=3\cdot54+36,\quad54=36+18,\quad36=2\cdot18.
\]
Отсюда \(18=54-36=54-(198-3\cdot54)=4\cdot54-198=4\cdot252-5\cdot198\). Деления находят НОД; обратная подстановка одновременно находит коэффициенты Безу. Такой способ называют \term{расширенным алгоритмом Евклида} (extended Euclidean algorithm).

\section*{Когда по модулю можно делить}
\begin{definition}
Целое \(u\) — \term{обратное к \(a\) по модулю \(m\)} (modular multiplicative inverse), если \(au\equiv1\pmod m\), где \(m\ge2\). Обратный класс, если он существует, единственен; разные целые представители отличаются на кратное \(m\).
\end{definition}

\begin{theorem}[Критерий обратимости]
При \(m\ge2\) число \(a\) имеет обратное по модулю \(m\) тогда и только тогда, когда \(\gcd(|a|,m)=1\).
\end{theorem}
\begin{proof}
Если \(au\equiv1\pmod m\), то \(au-mv=1\) для некоторого целого \(v\). Любой общий делитель \(a,m\) делит левую часть, а значит, и \(1\); НОД равен \(1\). Обратно, при НОД \(1\) тождество Безу даёт \(au+mv=1\), откуда \(au\equiv1\pmod m\).
Если \(u\) и \(w\) оба обратны, то \(u\equiv u(aw)=(ua)w\equiv w\pmod m\). Здесь применены только умножение и правила сравнений.
\end{proof}
Например, \(3\cdot5=15\equiv1\pmod7\). Поэтому \(3x\equiv4\pmod7\) можно умножить на \(5\), получив \(x\equiv6\pmod7\). У числа \(6\) нет обратного по модулю \(9\), поскольку НОД равен \(3\).

\begin{theorem}[Сокращение сравнения]
Если \(ac\equiv bc\pmod m\) и \(\gcd(|c|,m)=1\), то \(a\equiv b\pmod m\).
\end{theorem}
\begin{proof}
При \(m=1\) все целые числа сравнимы, и вывод очевиден. Пусть далее \(m\ge2\). Умножим сравнение на обратное к \(c\) по модулю \(m\), существующее по критерию. Получим \(a(cu)\equiv b(cu)\), а \(cu\equiv1\); значит, \(a\equiv b\).
\end{proof}
Без условия на НОД сокращение может изменить множество решений. Из \(2x\equiv2\pmod6\) следует \(x\equiv1\pmod3\); в классах по модулю \(6\) это два ответа, \(x\equiv1\) и \(x\equiv4\). Оставить только \(x\equiv1\pmod6\) означало бы потерять половину решений.

\section*{Уравнения в целых числах}
\begin{definition}
\term{Линейное диофантово уравнение} (linear Diophantine equation) с двумя неизвестными имеет вид \(ax+by=c\), где коэффициенты целые и требуется найти все целочисленные пары \((x,y)\).
\end{definition}

\begin{theorem}[Решения линейного диофантова уравнения]
Пусть \(a,b\) не равны нулю одновременно и \(d=\gcd(|a|,|b|)\). Уравнение \(ax+by=c\) разрешимо в целых числах тогда и только тогда, когда \(d\mid c\). Если \((x_0,y_0)\) — одно решение, то все решения задаются формулами
\[
x=x_0+\frac bd\,t,\qquad y=y_0-\frac ad\,t,\qquad t\in\mathbb Z.
\]
\end{theorem}
\par\addvspace{5pt}\begin{lrbox}{\CatFlowBox}\begin{minipage}[t]{\linewidth}
\setlength{\parskip}{5pt}\setlength{\columnsep}{4.5mm}\setlength{\intextsep}{0pt}
\begin{wrapfigure}[9]{r}{54mm}
\includegraphics[width=54mm,height=58mm,keepaspectratio]{images/unique/art-71.png}
\end{wrapfigure}
\textit{Доказательство.} Любая комбинация \(ax+by\) делится на \(d\), поэтому условие необходимо. Если \(c=dk\), умножение тождества Безу на \(k\) даёт решение.
Для описания всех решений сначала пусть \(a,b\ne0\). Вычитание двух уравнений даёт \((a/d)(x-x_0)=-(b/d)(y-y_0)\). Выполнено \(\gcd(|a|/d,|b|/d)=1\). Из тождества Безу следует: если \(v\ne0\), \(\gcd(|u|,|v|)=1\) и \(v\mid uz\), то \(v\mid z\). Поэтому \(x-x_0=(b/d)t\), а подстановка даёт \(y-y_0=-(a/d)t\). Обратная подстановка проверяет все такие пары. Если один коэффициент нулевой, неизвестная при нём свободна, а другая фиксирована; те же формулы перечисляют все решения, поскольку ненулевое отношение коэффициента к \(d\) равно \(1\) или \(-1\).\hfill\(\blacksquare\)
\par\WFclear\end{minipage}\end{lrbox}%
\ifdim\dimexpr\ht\CatFlowBox+\dp\CatFlowBox\relax<64mm\dp\CatFlowBox=\dimexpr64mm-\ht\CatFlowBox\relax\fi%
\noindent\usebox{\CatFlowBox}\par\addvspace{5pt}

Например, \(14\cdot2+21\cdot(-1)=7\). Все целые решения \(14x+21y=7\) имеют вид \(x=2+3t\), \(y=-1-2t\). Если по смыслу задачи нужны неотрицательные количества, ограничения \(x\ge0\), \(y\ge0\) накладывают \emph{после} нахождения общих целочисленных решений.

\Needspace{7\baselineskip}\section*{Задачи}
\begingroup\setlength{\parskip}{3pt}
\noindent\begin{minipage}[t]{\linewidth}\textbf{21.1.} Раздели \(47\) на \(5\) с остатком; укажи \(q,r\).\par\end{minipage}\par
\noindent\begin{minipage}[t]{\linewidth}\textbf{21.2.} Раздели \(-47\) на \(5\) с неотрицательным остатком.\par\end{minipage}\par
\noindent\begin{minipage}[t]{\linewidth}\textbf{21.3.} Найди остаток \(-3\) при делении на \(7\).\par\end{minipage}\par
\noindent\begin{minipage}[t]{\linewidth}\textbf{21.4.} Проверь сравнения \(22\equiv4\pmod6\) и \(22\equiv5\pmod6\).\par\end{minipage}\par
\noindent\begin{minipage}[t]{\linewidth}\textbf{21.5.} Найди остаток \(137+284\) по модулю \(9\), сначала заменив слагаемые их остатками.\par\end{minipage}\par
\noindent\begin{minipage}[t]{\linewidth}\textbf{21.6.} Найди остаток \(10^{20}\) по модулю \(9\).\par\end{minipage}\par
\noindent\begin{minipage}[t]{\linewidth}\textbf{21.7.} Найди остаток \(2^{10}\) по модулю \(7\), используя \(2^3\equiv1\pmod7\).\par\end{minipage}\par
\noindent\begin{minipage}[t]{\linewidth}\textbf{21.8.} Определи последнюю десятичную цифру числа \(7^{100}\).\par\end{minipage}\par
\par\addvspace{5pt}\begin{lrbox}{\CatFlowBox}\begin{minipage}[t]{\linewidth}
\setlength{\parskip}{5pt}\setlength{\columnsep}{4.5mm}\setlength{\intextsep}{0pt}
\begin{wrapfigure}[12]{r}{54mm}
\includegraphics[width=54mm,height=58mm,keepaspectratio]{images/unique/art-61.png}
\end{wrapfigure}
\textbf{21.9.} На 24-часовых часах сейчас \(23{:}00\). Какое время будет через \(8\) часов?\par

\textbf{21.10.} Узор повторяет цвета «синий, кремовый, охра, бирюзовый». Первый элемент синий. Какого цвета \(31\)-й элемент?\par

\textbf{21.11.} Сегодня понедельник. Какой день недели будет через \(100\) дней?\par

\textbf{21.12.} Найди наименьшее положительное \(x\), для которого \(x\equiv2\pmod3\) и \(x\equiv3\pmod5\).\par

\textbf{21.13.} Алгоритмом Евклида найди \(\gcd(252,198)\) и вырази его целочисленной линейной комбинацией исходных чисел.\par

\textbf{21.14.} Найди обратный класс к \(3\) по модулю \(7\).\par

\textbf{21.15.} Существует ли обратный к \(6\) по модулю \(9\)? Обоснуй.\par
\par\WFclear\end{minipage}\end{lrbox}%
\ifdim\dimexpr\ht\CatFlowBox+\dp\CatFlowBox\relax<64mm\dp\CatFlowBox=\dimexpr64mm-\ht\CatFlowBox\relax\fi%
\noindent\usebox{\CatFlowBox}\par\addvspace{5pt}

\noindent\begin{minipage}[t]{\linewidth}\textbf{21.16.} Реши \(3x\equiv4\pmod7\) во всех целых числах.\par\end{minipage}\par
\noindent\begin{minipage}[t]{\linewidth}\textbf{21.17.} Перечисли все классы решений \(2x\equiv2\pmod6\) по модулю \(6\).\par\end{minipage}\par
\noindent\begin{minipage}[t]{\linewidth}\textbf{21.18.} Реши \(2x\equiv1\pmod6\).\par\end{minipage}\par
\noindent\begin{minipage}[t]{\linewidth}\textbf{21.19.}\starmark{} Найди все целочисленные решения \(14x+21y=7\).\par\end{minipage}\par
\noindent\begin{minipage}[t]{\linewidth}\textbf{21.20.} Имеет ли целочисленные решения уравнение \(6x+10y=7\)?\par\end{minipage}\par
\noindent\begin{minipage}[t]{\linewidth}\textbf{21.21.}\starmark{} Наборы стоят \(12\) и \(18\) условных единиц. Нужно набрать ровно \(120\), покупая целые неотрицательные количества \(x,y\). Найди все пары количеств.\par\end{minipage}\par
\noindent\begin{minipage}[t]{\linewidth}\textbf{21.22.} Докажи признак делимости на \(3\) через сравнение \(10\equiv1\pmod3\).\par\end{minipage}\par
\noindent\begin{minipage}[t]{\linewidth}\textbf{21.23.} Докажи, что квадрат целого числа имеет по модулю \(4\) только остаток \(0\) или \(1\).\par\end{minipage}\par
\noindent\begin{minipage}[t]{\linewidth}\textbf{21.24.} Существуют ли целые \(x\), для которых \(x^2\equiv2\pmod4\)?\par\end{minipage}\par
\noindent\begin{minipage}[t]{\linewidth}\textbf{21.25.} Докажи, что из \(ac\equiv bc\pmod m\) можно сократить \(c\), если \(\gcd(|c|,m)=1\). Дай пример потери решения без этого условия.\par\end{minipage}\par
\noindent\begin{minipage}[t]{\linewidth}\textbf{21.26.}\starmark{} Найди все классы решений \(x^2\equiv1\pmod8\), обосновав ответ без перебора больших чисел.\par\end{minipage}\par
\noindent\begin{minipage}[t]{\linewidth}\textbf{21.27.}\starmark{} Найди все неотрицательные целочисленные решения \(3x+5y=31\).\par\end{minipage}\par
\noindent\begin{minipage}[t]{\linewidth}\textbf{21.28.}\starmark{} Для любых целых \(a,b\) докажи, что система \(x\equiv a\pmod3\), \(x\equiv b\pmod5\) имеет ровно один класс решений по модулю \(15\). Подсказка: проверь \(x=10a+6b\).\par\end{minipage}\par
\endgroup
